\subsubsection{正交多项式}\label{subsec:orthogonal-polynomials}

由 Weierstrass 逼近定理, 多项式在 $L^2(-1,1)$ 中稠密. 因而自然可以对序列 $1,x,x^2,\ldots$ 应用 Schmidt 正交化过程, 得到由多项式组成的一组基. 我们断言, 除去因子, 这些正交多项式由下式给出
\begin{equation}\label{eq:legendre-polynomial-rodrigues-formula}
P_n(x)=\dfrac1{2^nn!}\dfrac{\mathrm{d}^n(x^2-1)^n}{\mathrm{d}x^n}.
\end{equation}


首先, 显然 $P_n$ 为 $n$ 次多项式. 此外, 分部积分表明
\begin{equation}\label{eq:legendre-lower-monomial-orthogonality}
\int_{-1}^1 P_n(x)x^m\,\mathrm{d}x=0
\end{equation}
对每个 $m<n$ 成立; 因而还有
\begin{equation}\label{eq:legendre-polynomial-pairwise-orthogonality}
\int_{-1}^1 P_n(x)P_m(x)\,\mathrm{d}x=0
\end{equation}
对 $m<n$ 成立, 即 $P_n$ 在 $L^2(-1,1)$ 中正交. $P_n$ 称为 Legendre 多项式. 最前面的几个为
\begin{equation}\label{eq:first-five-legendre-polynomials}
\begin{gathered}
P_0(x)=1,\quad P_1(x)=x,\quad P_2(x)=\dfrac32x^2-\dfrac12,\\
P_3(x)=\dfrac52x^3-\dfrac32x,\quad P_4(x)=\dfrac{35}8x^4-\dfrac{15}4x^2+\dfrac38.
\end{gathered}
\end{equation}


Legendre 多项式尚未归一化. 反复分部积分可得
\begin{equation}\label{eq:legendre-square-integral-reduction}
\int_{-1}^1 P_n^2(x)\,\mathrm{d}x=\dfrac{(2n)!}{2^{2n}(n!)^2}\int_{-1}^1(1-x^2)^n\,\mathrm{d}x.
\end{equation}
注意 $(1-x^2)^n=(1-x)^n(1+x)^n$, 再反复分部积分, 即可求出 $(1-x^2)^n$ 的积分. 最终结果为
\begin{equation}\label{eq:legendre-square-integral-value}
\int_{-1}^1 P_n^2(x)\,\mathrm{d}x=\dfrac2{2n+1}.
\end{equation}


其他许多正交多项式对应用也很重要. 这些多项式在加权 $L^2$ 空间中正交.

\begin{Definition}
\label{def:weighted-lp-two-space}
设 $\Omega$ 为 $\mathbb{R}^m$ 中的开集, $w$ 为从 $\Omega$ 到 $\mathbb{R}^+$ 的连续函数. 定义
\begin{equation}\label{eq:weighted-continuous-lp-two-space}
\widetilde L_w^2(\Omega):=\left\{u\in C(\overline\Omega)\mid\int_\Omega w(\bm{x})|u(\bm{x})|^2\,\mathrm{d}\bm{x}<\infty\right\}.
\end{equation}
内积定义为
\begin{equation}\label{eq:weighted-lp-two-inner-product}
(u,v)=\int_\Omega w(\bm{x})\overline{u(\bm{x})}v(\bm{x})\,\mathrm{d}\bm{x}.
\end{equation}
空间 $L_w^2(\Omega)$ 是 $\widetilde L_w^2(\Omega)$ 的完备化.
\end{Definition}

对任意权函数 $w$ 和任意区间 $(a,b)$, 现在都可以通过对序列 $1,x,x^2,\ldots$ 正交化来定义正交多项式 (当然前提是 $w$ 使多项式属于 $L_w^2(\Omega)$). 以下情形尤其重要:
